% ============================================================
% Appendix
% ============================================================

\section{DATA COLLECTION SOURCE CODE}

\subsection{Main Collection Script (3-collect\_dataset\_v2.py)}

The complete source code for the multi-modal data collection pipeline is provided below. This script implements the optimized data collection system described in Chapter III.

\begin{lstlisting}[language=Python, caption={Collection Configuration Class}]
@dataclass
class CollectionConfig:
    """Collection configuration dataclass."""
    vessel_name: str = "Vessel1"
    front_cam_idx: int = 0
    lidar_name: str = "top_lidar"
    num_frames: int = 200
    move_seconds: float = 0.3
    thrust: float = 0.3
    angle: float = 0.6
    trajectory_mode: str = "circle"
    cam_w: int = 640
    cam_h: int = 480
    cam_fov: float = 90.0
    dataset_root: str = "dataset"
    max_retries: int = 3
    retry_delay: float = 1.0
    max_fail_ratio: float = 0.2
\end{lstlisting}

\begin{lstlisting}[language=Python, caption={Camera Intrinsics Computation}]
def compute_intrinsics(w: int, h: int, fov_deg: float) -> Dict:
    """
    Compute camera intrinsics matrix K.
    Formula: fx = w / (2 * tan(FOV/2))
    """
    fov_rad = math.radians(fov_deg)
    fx = w / (2 * math.tan(fov_rad / 2))
    return {
        "fx": fx, "fy": fx,
        "cx": w / 2.0, "cy": h / 2.0,
        "width": w, "height": h,
        "fov_deg": fov_deg,
    }
\end{lstlisting}

\begin{lstlisting}[language=Python, caption={Image Decoding with Color Space Conversion}]
def decode_image_safe(response, name: str = "") -> np.ndarray:
    """
    Safely decode images, handling RGB/RGBA formats
    and converting to BGR for OpenCV compatibility.
    """
    raw = np.frombuffer(response.image_data_uint8, dtype=np.uint8)
    h, w = response.height, response.width
    total = len(raw)

    expected_3ch = h * w * 3
    expected_4ch = h * w * 4

    if total == expected_4ch:
        # RGBA - convert to BGR
        img = raw.reshape(h, w, 4)[:, :, :3]
        img = cv2.cvtColor(img, cv2.COLOR_RGB2BGR)
    elif total == expected_3ch:
        # RGB - convert to BGR
        img = raw.reshape(h, w, 3)
        img = cv2.cvtColor(img, cv2.COLOR_RGB2BGR)
    else:
        # Unknown format - attempt recovery
        channels = max(1, total // (h * w))
        img = raw.reshape(h, w, channels)
        if channels > 3:
            img = img[:, :, :3]
        if len(img.shape) == 3 and img.shape[2] == 3:
            img = cv2.cvtColor(img, cv2.COLOR_RGB2BGR)
    return img
\end{lstlisting}

\subsection{ASVSim Settings Configuration}

The complete \texttt{settings.json} configuration for optimized data collection:

\begin{lstlisting}[language={}, caption={Optimized ASVSim Settings}]
{
  "SettingsVersion": 2.0,
  "SimMode": "Vessel",
  "PhysicsEngineName": "VesselEngine",
  "ViewMode": "SpringArmChase",
  "InitialInstanceSegmentation": true,
  "ClockSpeed": 1,
  "Wind": { "X": 0, "Y": 0, "Z": 0 },

  "PawnPaths": {
    "BlueResearchBoat": {
      "PawnBP": "Class'/AirSim/Blueprints/BP_VesselPawn.BP_VesselPawn_C'"
    }
  },

  "Vehicles": {
    "Vessel1": {
      "VehicleType": "MilliAmpere",
      "AutoCreate": true,
      "PawnPath": "BlueResearchBoat",
      "HydroDynamics": { "hydrodynamics_engine": "FossenCurrent" },
      "X": 0, "Y": 0, "Z": 0,
      "Sensors": {
        "top_lidar": {
          "SensorType": 6, "Enabled": true,
          "NumberOfLasers": 16,
          "PointsPerSecond": 300000,
          "RotationsPerSecond": 10,
          "Range": 100.0,
          "DrawDebugPoints": false,
          "DataFrame": "SensorLocalFrame",
          "X": 0.0, "Y": 0.0, "Z": -1.0
        },
        "front_camera": {
          "SensorType": 1, "Enabled": true,
          "Width": 640, "Height": 480,
          "X": 3.5, "Y": 0.0, "Z": -2.5,
          "Pitch": 0.0, "Roll": 0.0, "Yaw": 0.0
        }
      }
    }
  },

  "CameraDefaults": {
    "CaptureSettings": [
      {
        "ImageType": 0,
        "Width": 640, "Height": 480,
        "FOV_Degrees": 70,
        "LumenGIEnable": false,
        "LumenReflectionEnable": false,
        "MotionBlurAmount": 0,
        "ChromaticAberrationScale": 0.0,
        "AutoExposureBias": 0.0
      },
      {
        "ImageType": 1,
        "Width": 640, "Height": 480,
        "FOV_Degrees": 70,
        "LumenGIEnable": false,
        "LumenReflectionEnable": false,
        "MotionBlurAmount": 0
      },
      {
        "ImageType": 5,
        "Width": 640, "Height": 480,
        "FOV_Degrees": 70,
        "LumenGIEnable": false,
        "LumenReflectionEnable": false
      }
    ]
  }
}
\end{lstlisting}

\subsection{Dataset Validation Script}

The validation script checks dataset integrity:

\begin{lstlisting}[language=Python, caption={Dataset Validation Function}]
def validate_dataset(dirs: Dict[str, Path],
                     expected_frames: int,
                     collected_frames: int) -> Dict:
    """Validate dataset completeness and quality."""
    results = {
        "valid": True,
        "issues": [],
        "stats": {},
    }

    # Check file counts
    rgb_files = sorted(list(dirs["rgb"].glob("*.png")))
    depth_files = sorted(list(dirs["depth"].glob("*.npy")))

    results["stats"]["rgb_count"] = len(rgb_files)
    results["stats"]["depth_count"] = len(depth_files)

    # Validate image quality
    if rgb_files:
        img = cv2.imread(str(rgb_files[0]))
        if img.shape != (480, 640, 3):
            results["valid"] = False
            results["issues"].append(f"Wrong dimensions: {img.shape}")

    # Check pose file
    poses_path = dirs["base"] / "poses.json"
    if poses_path.exists():
        with open(poses_path, "r") as f:
            poses = json.load(f)
        if len(poses) != collected_frames:
            results["valid"] = False
            results["issues"].append("Pose count mismatch")

    return results
\end{lstlisting}
